
\section{Proofs for overparameterized least-squares}\label{sec proof thm 1}
In this section, we assume the linear Gaussian problem (LGP) of Definition \ref{def LGP}, the overparameterized regime $k=p>n$ and the min-norm model $\bth$ of \eqref{eq:min_norm}. We prove Theorem \ref{thm:master_W2} that derives the asymptotic DC of $\bth$ and we show how this leads to sharp formulae for the risk of the Magnitude- and Hessian-pruned models.


\subsection{Notation and Assumptions}\label{sec:ass_app}
%\noindent\textbf{Notation.}~
For the reader's convenience, we recall some necessary notation and assumptions from Section \ref{sec main}.  We say that a function $f:\R^m\rightarrow\R$ is pseudo-Lipschitz of order $k$, denoted $f\in\rm{PL}(k)$, if there is a constant $L>0$ such that for all $\x,\y\in\R^m$,
$
|f(\x) - f(\y)|\leq L(1+\tn{\x}^{k-1}+\tn{\y}^{k-1})\|\x-\y\|_2 
$
(See also Section \ref{SM useful fact}). 
We say that a sequence of probability distributions $\nu_p$ on $\R^m$ {converges in $W_k$} to 
$\nu$, written $\nu_p\stackrel{W_k}{\Longrightarrow} \nu$, if $W_k(\nu_p,\nu) \rightarrow 0$ as $p \rightarrow \infty$. An equivalent definition is that, for any $f\in\rm{PL}(k)$, $\lim_{p\rightarrow\infty}\E f(X_p)=\E f(X)$, where expectation is with respect to $X_p\sim\nu_p$ and $X\sim\nu$ (e.g., \cite{montanari2017estimation}).  
%
Finally, recall that a sequence of probability distributions $\nu_n$ on $\R^m$ \emph{converges weakly} to $\nu$, if for any bounded Lipschitz function $f$: $\lim_{p\rightarrow\infty}\E f(X_p)=\E f(X)$, where expectation is with respect to $X_p\sim\nu_p$ and $X\sim\nu$.
%
Throughout, we use $C,C',c,c'$ to denote absolute constants (not depending on $n,p$) whose value might change from line to line. 


We focus on a double asymptotic regime where:
$$n,p,s\rightarrow\infty \text{ at fixed overparameterization ratio } \kappa:=p/n>1 \text{ and sparsity level } \alpha:=s/p\in(0,1).$$  
%
For a sequence of random variables $\mathcal{X}_{p}$ that converge in probability to some constant $c$ in the limit of Assumption \ref{ass:linear} below, we  write $\mathcal{X}_{p}\rP c$. For a sequence of event $\Ec_p$ for which $\lim_{p\rightarrow}\Pro(\Ec_p) = 1$, we say that $\Ec_p$ occurs \emph{with probability approaching 1}. For this, we will often use the shorthand ``wpa 1". 

\vspace{5pt}
Next, we recall the set of assumption under which our analysis applies:


\asstwo*
\assthree*


%
%\vp
%\textbf{(A1)[Diagonal covariance].}~~The covariance matrix $\Sigmab$ is diagonal.
%
%%
%\vp
%\textbf{\ref{ass:mu}[Boundedness and empirical distribution].}~~There exist constants $\Sigma_{\min},\Sigma_{\max}\in(0,\infty)$ such that:
%$
%\Sigma_{\min}\leq\bSi_{i,i}\leq \Sigma_{\max}.
%$
%for all $i\in[p].$ 
%%There exists absolute constant  $C>0$ such that $\sqrt{p}|\betas_i|<C$, for all $i\in[p]$. 
%\cts{Furthermore, the joint empirical distribution of $\{(\lai,\sqrt{p}\betas_i)\}_{i\in[p]}$ converges weakly to a probability distribution $\mu$ on $\R_{>0}\times\R$ with bounded $4$th moment 
%%$\E_{(\Lambda,B)}\left[\left(\sqrt{\Lambda^2+B^2}\right)^3\right]$, 
%and assume that as $p\rightarrow\infty$, the $4$th moment of the empirical distribution converges to the $4$th moment of $\mu$, i.e.,
%$$
%\frac{1}{p}\sum_{i=1}^p\left(\sqrt{(\sqrt{p}\betas_i)^{2}+\lai^2}\right)^4 \rP \E_{(\Lambda,B)\sim\mu}\left[\left(\sqrt{\Lambda^2+B^2}\right)^4\right]<\infty.
%$$
%}
% for some $k\geq 3$.}

%\cts{Furthermore, the joint empirical distribution of $\{(\lai,\sqrt{p}\betas_i)\}_{i\in[p]}$ converges weakly to a probability distribution $\mu$ on $\R_{>0}\times\R$ with bounded $(2k-2)$-th moment, and assume that 
%$$
%\frac{1}{p}\sum_{i\in[p]}\sum_{i=1}^p(\sqrt{p}\betas_i)^{2k-2} \rP \E_{(\Lambda,B)\sim\mu}\left[ B^{2k-2}\right]<\infty,
%$$
%as $p\rightarrow\infty$ for some $k\geq 3$.}


\noindent We remark that Assumption \ref{ass:mu} above implies (see \cite[Lem.~4]{bayati2011dynamics} and \cite[Lem.~A3]{javanmard2013state}) that for any pseudo-Lipschitz function $\psi:\R^2\rightarrow\R$ of order $4$, i.e., $\psi\in\rm{PL}(4)$:
$$
\frac{1}{p}\sum_{i=1}^p\psi(\lai,\sqrt{p}\betas_i) \rP \E_{(\Lambda,B)\sim\mu}\left[ \psi(\Lambda,B)\right].
$$




\subsection{Asymptotic distribution and risk characterizations}
%, \fx{which is a special case of Theorem \ref{thm:master_W2}, tailored to our specific need in this paper, that is, characterizing the risk of pruned solutions.} \somm{Theorem 1 should be the special case. Not the other way around!!!}

{In this section, we prove our main result Theorem \ref{thm:master_W2}. Recall that $\bth$ is the min-norm solution. Since the distribution of $\bth$ depends on the problem dimensions (as it is a function of $\X,\y$), when necessary, we will use $\bth_n$ notation to make its dimension dependence explicit.} Let $\bth^P=\Pc(\bth)$ be a pruned version of the min-norm solution $\bth$. Recall from Section \ref{sec:risk}, that the first crucial step in characterizing the risk $\Lc(\bth^P)$ is studying the risk $\Lc(\bth^\Tc_t)$ of a threshold-based pruned vector. 
%Theorem \ref{thm:master_W2} below shows how this is possible. 

To keep things slightly more general,  consider $\bth^{g}$ defined such that $\sqrt{p}\bth^{g}=g(\sqrt{p}\bth)$, where $g$ is a Lipschitz function acting entry-wise on $\bth$ (for example, $g$ can be the {(arbitrarily close Lipschitz approximation of the)} thresholding operator $\Tc_t$ of Section \ref{sec:risk}). Then,  the risk of $\bth^g$ can be written as
\begin{align}
\Lc(\bth^g) &= \E_{\Dc}[(\x^T(\bt^\star-\bth^g) + \sigma z)^2] = \sigma^2 + (\bt^\star-\bth^g)^T\Sigmab(\bt^\star-\bth^g) \nn \\
  &= \sigma^2 + \frac{1}{p}\sum_{i=1}^p\Sigmab_{i,i}\big(\sqrt{p}\betas_i-g(\sqrt{p}\bth_i)\big)^2\nn
  \\
    &=: \sigma^2 + \frac{1}{p}\sum_{i=1}^p f\big(\sqrt{p}\bth_i,\sqrt{p}\betas_i,\Sigmab_{i,i}\big),\label{eq:risk_app_f}
% \\
%  &\rP \sigma^2 + \E\left[ \Lambda\left(B-\Tc_t(X_{\kappa,\sigma^2})\right) \right]\,.
\end{align}
where in the last line, we defined the function $f$ as $f=\fl\in \Fl\subset\Plt$ given by
\[
\fl(x,y,z) := z(y-g(x))^2\quad\text{where}~g~\text{is Lipschitz.}
\]
Here, recall the definition of the families $\Fl$ in \eqref{eq:fdef} and $\Plt$ in \eqref{eq:pdef}.

The following theorem establishes the asymptotic limit of \eqref{eq:risk_app_f}. For the reader's convenience, we repeat the notation introduced in Definition \ref{def:Xi}.
%\begin{definition}[Asymptotic DC -- Overparameterized regime]\label{def:Xi}
Let random variables $(\Lambda,B)\sim \mu$ (where $\mu$ is defined in Assumption \ref{ass:mu}) and fix $\kappa>1$. Define parameter $\xi$ as the unique positive solution to the following equation
%\begin{align}\label{eq:ksi}
$$
\E_{\mu}\Big[ \big({1+(\xi\cdot\Lambda)^{-1}}\big)^{-1} \Big] = {\kappa^{-1}}\,.
$$
Further define the positive parameter $\gamma$ as follows:
$$
\hspace{-0.1in}\gamma := 
%\frac{\sigma^2 + \E_{\mu}\left[\frac{N^2}{(\Lambda^{-1}+\xi)^2}\right]}{1-\kappa\E_{\mu}\left[\frac{1}{\left(1+(\xi\Lambda)^{-1}\right)^2}\right]}  
\Big({\sigma^2 + \E_{\mu}\Big[\frac{B^2\Lambda}{(1+\xi\Lambda)^2}\Big]}\Big)\Big/\Big({1-\kappa\E_{\mu}\Big[\frac{1}{\left(1+(\xi\Lambda)^{-1}\right)^2}\Big]}\Big).
$$
%\ct{@Samet: In your notation $\gamma$ should be $\Gamma\kappa$. Everything matches except the second term in the numerator. In your notation, that second term becomes $\int_0^1\bz^2(t)\bt^2(t) \Sigmab(t) dt$ (note the extra $\Sigmab(t)$ compared to your formula)
%}
With these and $H\sim\Nn(0,1)$, define the random variable
$$
X_{\kappa,\sigma^2}:=X_{\kappa,\sigma^2}(\Lambda,B,H) := \Big(1-\frac{1}{1+ \xi\Lambda}\Big) B + \sqrt{\kappa}\frac{\sqrt{\gamma}\,\Lambda^{-1/2}}{1+(\xi\Lambda)^{-1}} H, 
$$
and let $\Pi_{\kappa,\sigma^2}$ be its distribution.
%\end{definition}


%Recall the definition of the measure $\Pi_{\kappa,\sigma^2}$ in Definition \ref{def:Xi}. Then, 
%$\hat\Pi_n(\y,\X,\betas,\Sigmab)$ converges in Wasserstein-2 distance to $\Pi_{\kappa,\sigma^2}\otimes\mu$.
%Let $\cal{H}$ be the hard-thresholding operation. Draw $\hat{\bt}\sim \Dca(\bt,\bSi)$.
%\[
%\text{risk}_{\bSi^p}(\Hc(\hat{\bt}^p))\rightarrow \text{risk}_{\bSi}(\Hc(\hat{\bt}))
%\]
%\end{theorem}
%
%\begin{corollary} 
%Moreover, for the function $\fl:\R^3\rightarrow\R$, defined in \eqref{eq:fdef} it holds that
%Specifically, for any function $f:\R^k\rightarrow\R$, $f\in\rm{PL}(k)$, it holds that
%where the expectation is over $(\Lambda,B,H)\sim\mu\otimes\Nn(0,1).$
%Specifically, letting $\cal{H}:\R\rightarrow\R$ denote the hard-thresholding operation, it holds that
%\begin{align}
%\| \cal{H}( \hat\w^p )\|_2^2 \rP %\E\left[\left(\cal{H}\left(X_{\kappa,\sigma^2}(\Lambda,N,H) \right)\right)^2\right]
%\end{align}
%%\end{corollary}

\mainthm*

Before we prove the theorem, let us show how it immediately leads to a sharp prediction of the risk behavior. Indeed, a direct application of \eqref{eq:thm} for $f=\fl$ to \eqref{eq:risk_app_f} shows that
\begin{align}
\Lc(\bth^g)\rP \sigma^2 +  \E_{(\Lambda,B,H)\sim\mu\otimes\Nn(0,1)}\left[\fl(X_{\kappa,\sigma^2},B,\Lambda) \right] = \sigma^2 +  \E_{(\Lambda,B,H)\sim\mu\otimes\Nn(0,1)}\left[\Sigma\left(B-g(X_{\kappa,\sigma^2})\right)^2 \right].\label{eq:risk_app_f2}
\end{align}

We further remark on the following two consequences of Theorem \ref{thm:master_W2}. 

First, since \eqref{eq:thm} holds for any $\rm{PL}(2)$ function, we have essentially shown that $\hat\Pi_n(\y,\X,\betas,\Sigmab)$ converges in Wasserstein-2 distance to $\Pi_{\kappa,\sigma^2}\otimes\mu$, where recall that $\Pi_{\kappa,\sigma^2}$ is the distribution of the random variable $X_{\kappa,\sigma^2}$.

Second, the theorem implies that the empirical distribution of $\sqrt{p}\bth_n$ converges weakly to $\Pi_{\kappa,\sigma^2}$. To see this, apply \eqref{eq:thm} for the $\rm{PL}(2)$ function $f(x,y,z) = \psi(x)$ where  $\psi:\R\rightarrow\R$ is a bounded Lipschitz test function. 

\subsection{Proof of Theorem \ref{thm:master_W2}}
\vspace{5pt}
Let $\X\in\R^{n\times p}$ have zero-mean and normally distributed rows with a diagonal covariance matrix $\bSi=\E[\x\x^T]$. Given a ground-truth vector $\betas$ and labels $\y=\X\betas+\sigma \z,~\z\sim\Nn(0,\Iden_n)$, we consider the least-squares problem subject to the minimum Euclidian norm constraint (as $\kappa=p/n>1$) given by
\begin{align}\label{eq:PO_beta}
\min_{\bt}\frac{1}{2}\tn{\bt}^2\quad\text{subject to}\quad \y=\X\bt.
\end{align}
%, we solve the least-squares problem  \eqref{optim me2} subject to the . 
It is more convenient to work with the following change of variable: 
\begin{align}\label{eq:w}
\w:=\sqrt{\Sigmab}(\bt-\betas).
\end{align}
 With this, the optimization problem  in \eqref{eq:min_norm} can be rewritten as
\begin{align}\label{eq:PO}
\Phi(\X)=\min_{\w} \frac{1}{2}\tn{\Sigmab^{-1/2}\w+\betas}^2\quad\text{subject to}\quad \Xb\w=\sigma \z,
\end{align}
where we set $\Xb=\X\Sigmab^{-1/2}\distas\Nn(0,1)$. First, using standard arguments, we show that the solution of \eqref{eq:PO} is bounded. Hence, we can constraint the optimization in a sufficiently large compact set without loss of generality.


\begin{lemma}[Boundedness of the solution]\label{lem:bd_PO}
Let $\wh_n:=\wh_n(\X,\z)$ be the minimizer in \eqref{eq:PO}. Then, with probability approaching 1, it holds that $\wh_n\in\Bc$, where
$$\Bc:=\left\{\w\,|\,\|\w\|_2\leq B_{+} \right\},\qquad B_+:=5\sqrt{\frac{\Sigma_{\max}}{\Sigma_{\min}}}\frac{\sqrt{\kappa}+1}{\sqrt{\kappa}-1}(\sqrt{\Sigma_{\max}\E\left[B^2\right]} + \sigma).
% + \sqrt{\Sigma_{\max}\E\left[B^2\right]}.
$$ 
\end{lemma}
\begin{proof}
First, we show that the min-norm solution $\bth=\X^T(\X\X^T)^{-1}\y$ of \eqref{eq:PO_beta} is bounded. Note that $\kappa>1$, thus $\X\X^T$ is invertible wpa 1. We
have,
\begin{align}
\tn{\bth_n}^2 = \y^T(\X\X^T)^{-1}\y \leq \frac{\tn{\y}^2}{\la_{\min}(\X\X^T)}  = \frac{\tn{\y}^2}{\la_{\min}(\Xb\Sigmab\Xb^T)} \leq  \frac{\tn{\y}^2}{\la_{\min}(\Xb\Xb^T)\,\Sigma_{\min}} = \frac{\tn{\y}^2}{\sigma_{\min}^2(\Xb)\,\Sigma_{\min}}. \label{eq:Ubb}
\end{align}
But, wpa 1, 
$
\sigma_{\min}(\Xb)/\sqrt{n} \geq \frac{1}{2}\left(\sqrt{\kappa}-1\right).
$
Furthermore, 
$
\|\y\|_2 \leq \|\Xb\Sigmab^{1/2}\betas\|_2 + \sigma\|\z\|_2 \leq \sigma_{\max}(\Xb)\sqrt{\Sigma_{\max}} \|\betas\|_2 + \sigma\|\z\|_2.
$
Hence, wpa 1,
$$
\|\y\|_2/\sqrt{n} \leq 2(\sqrt{\kappa}+1)\sqrt{\Sigma_{\max}} \sqrt{\E\left[B^2\right]} + 2\sigma,
$$
where we used the facts that wpa 1: $\|z\|_2/\sqrt{n}\rP 1$, $\sigma_{\max}(\Xb)<\sqrt{2n}(\sqrt{\kappa}+1)$ and \cts{by Assumption \ref{ass:mu}}:
$$
\|\betas\|_2^2 = \frac{1}{p}\sum_{i=1}^{p}(\sqrt{p}\betas_i)^2 \rP \E\left[B^2\right].
$$
Put together in \eqref{eq:Ubb}, shows that
\begin{align}
\tn{\bth_n} < \frac{2(\sqrt{\kappa}+1)\sqrt{\Sigma_{\max}} \sqrt{\E\left[B^2\right]} + 2\sigma}{\sqrt{\Sigma_{\min}}(\sqrt{\kappa}-1)/2} =: \tilde{B}_+.\label{eq:bd_beta}
\end{align}
Recalling that $\wh_n= \sqrt{\Sigmab}\bth_n - \sqrt{\Sigmab}\betas$, we conclude, as desired, that wpa 1,
$
\tn{\wh_n} \leq \sqrt{\Sigma_{\max}}\tilde{B}_+ + \sqrt{\Sigma_{\max}} \sqrt{\E[B^2]} \leq B_+.
$
%{\color{red}Missing proof}
\end{proof}

%%%%%%%%%%%%%%%%%%%%%%%%%%%%%%%%%
%%%% Samet's k-th bounded %%%%%%%%%%%%%%%%%
%%\begin{lemma}
%\end{lemma}
\so{\begin{lemma} \label{subexp bth}Let $\bth=\X^\dagger \y$ where $\y=\X\bts+\sigma \z$ where $\X,\z\distas\Nn(0,1)$. As $p,n\rightarrow \infty$, wpa.~1, there exists a constant $C>0$ such that, we have that
\[
\frac{1}{p}\sum_{i=1}^p|\sqrt{p}\bth_i|^k\leq C.
\]
\end{lemma}}
\begin{proof} We first show that entries of $\sqrt{p}\bth$ have finite $k$th moments for any $k\geq 0$. We split the proof in two components by writing $\bth=\ab+\sigma \bb$. 

First, let us focus on $\bb=\X^\dagger\z$. Write singular value decomposition $\X=\Ub\La\V^T$. Then $\bb=\V^T\La\Ub(\Ub\La^2\Ub^T)^{-1}\z=\V^T\La^{-1}\z'$ where $\z'=\Ub^T\z\distas\Nn(0,1)$. First, wpa.~1, $C\sqrt{p} \|\X\|\geq \sigma_{\min}(\X)\gtrsim c\sqrt{p}$ for some constants $C,c>0$ as the covariance $\bSi$ are bounded above and below by constants. To show subgaussianity, we fix unit length $\vb$ and consider the variable
\[
X=\vb^T\bb=\vb^T\V^T\La^{-1}\z'.
\]
by fixing all but $\z'$. Since $\z'$ is Gaussian, the subgaussian norm obeys 
\begin{align}
\tsub{\bb}=\sup_{\tn{\vb}=1}\tsub{X}&\lesssim \sup_{\tn{\vb}=1}\|\vb^T\V^T\La^{-1}\|_{\ell_\infty}\\
&\leq\sup_{\tn{\vb}=1} \tn{\vb^T\V^T}\sigma_{\min}(\X)^{-1}\\
&\leq \sigma_{\min}(\X)^{-1}\\
&\lesssim 1/\sqrt{p}.
\end{align}

To proceed, we address the subgaussianity of the main term 
\begin{align}
\ab&=\X^\dagger\X\bts=\X^T(\X\X^T)^{-1}\X\bts\\
&=\sqrt{\bSi}\Xb^T(\Xb\bSi\Xb^T)^{-1}\Xb\sqrt{\bSi}\bts\\
&=\sqrt{\bSi}\Qb^T (\Qb\bSi\Qb^T)^{-1}\Qb\sqrt{\bSi}\bts\\
&=\Vb^T\Vb\bts\\
&=\sum_{i=1}^n \vb_i\vb_i^T\bts.
\end{align} 
Here, $\Vb$ is the unitary matrix that is the range space of $\X$ and $\Qb$ is the unitary matrix that is the row space of $\Xb$. $(\vb_i)_{i=1}^n$ are rows of $\Vb$. Observe have that the range space of $\Vb$ is statistically identical to the range space of $\Qb\sqrt{\bSi}$. We will work with the properties of the row space $\Vb$ rather than the original design matrix $\X$. 

\noindent\textbf{Expectation of $\ab$:} We claim that $\E[\ab]=\m\odot \bts$ where $1\leq \m_j\geq 0$ for all $1\leq j\leq p$. Thus, by assumption on $\bts$, this would mean that $\E[|\sqrt{p}\E[\ab_i]|^k]$ and $p^{-1}\E[\|\sqrt{p}\E[\ab]\|_{\ell_k}^k]$ is bounded by a constant wpa.~1. To see this, observe that
\[
\E[\ab]=\sum_{i=1}^n\E[\vb_i\vb_i^T]\bts=\sum_{j=1}^p\sum_{i=1}^n\E[\vb_i\vb_i^T]\bts_j \eb_j,
\]
where $\eb_j$ is the $j$'th element of the standard basis. Note that the entries of $\vb_i$ have i.i.d.~Bernoulli signs same as the entries of the rows of $\Qb$. Thus $\E[\vb_i\vb_i^T]$ is a diagonal matrix $\Db^i$. This implies that, $\E[\X^\dagger\X]$ is diagonal and the $j$th index of $\bts$ is the only entry responsible from $\ab_j$. Thus
\[
\m_j=\sum_{i=1}^n \Db^i_{j,j}.
\]
Finally, $\sum_{i=1}^n \Db^i_{j,j}$ is equal to $\tn{\vb^\text{col}_j}^2\leq 1$ where $\vb^\text{col}_j$ is the $j$th column of $\Vb$ proving the result.

\noindent\textbf{Subexponentiality of $\ab$:} To conclude the proof, we show that $\ab$ is subexponential. Fix a unit length vector $\eb$ from standard basis. The claim is same as proving $\te{\sqrt{p}(X-\E[X])}\lesssim 1$ where
\begin{align}
X&=\eb^T\X^\dagger\X\bts=\eb^T\sqrt{\bSi}\Qb^T (\Qb\bSi\Qb^T)^{-1}\Qb\sqrt{\bSi}\bts.
\end{align}
\so{We will focus on the scenario $\bSi=\Iden_p$. General $\bSi$ is handled by the fact that $\bSi$ is upper and lower bounded by identity (thus isometric).} Thus, consider
\begin{align}
X&=\eb^T\X^\dagger\X\bts=\eb^T\Qb^T\Qb\bts\\
&=\frac{\tn{\Qb(\bts+\eb)}^2-\tn{\Qb(\bts-\eb)}^2}{4}.
\end{align}
Thus, our claim is same as showing
\[
\te{\tn{\Qb\vb}^2-\E[\tn{\Qb\vb}^2]}\lesssim 1/\sqrt{p}.
\]
for fixed $\vb=\bts+\vb$. Here, we can swap the randomness and assume fixed $\Qb$ and uniformly random $\vb$ entries of $\vb$ are equal to $\vb_i=\tn{\vb}\frac{\g_i}{\tn{\g}}$ for $\g\sim\Nn(0,\Iden_p)$ and $\tn{\vb}\lesssim 1$. Specifically, fixing $\Qb$ to be first $n$ elements of standard basis, we obtain
\[
\frac{1}{\tn{\vb}^2}(\tn{\Qb\vb}^2-\E[\tn{\Qb\vb}^2])=\frac{\sum_{i=1}^n \g_i^2}{\tn{\g}^2}-\frac{n}{p}=\frac{\sum_{i=1}^n \g_i^2-\frac{n}{p}\tn{\g}^2}{\tn{\g}^2}.
\]
To proceed, observe that both numerator and denominator are sub-exponential variables and concentrate accordingly to find that
\[
\Pro(|\tn{\Qb\vb}^2-\E[\tn{\Qb\vb}^2]|\gtrsim t/\sqrt{p})\leq \e^{-t}.
\]
Thus, entries of $\ab$ obeys $\te{X\sqrt{p}}\lesssim 1$ and together with bounded expectations, we obtain $\E[|\sqrt{p}X|^k]\lesssim 1$.

Finally, combining this with the $\bb$ term, we find the desired result $\E[|\sqrt{p}\bth_i|^k]\lesssim 1$.
%The final line follows from the fact that, since $\bSi_{i,i}$ are bounded below, wpa.~1, $\sigma_{\min}(\X)^{-1}\lesssim 1/\sqrt{p}$.

\so{To conclude with the proof, we use two things. First, $\frac{1}{p}\sum_{i=1}^p\bth_i$ weakly converges to the asymptotic limit given by Def.~\ref{def:Xi}. The proof is omitted for brevity but follows the exact same line of argument as the proof of Theorem \ref{thm:master_W2}. The main difference is that, for weak convergence proof, we don't require $k$th moments or pseudo-Lipschitz conditions (thus, the proof does not require this lemma and there is no circular logic). Once the weak convergence is established, we use Lemma 5 of \cite{bayati2011dynamics}. Since $k$th moments of $\bth_i$ are bounded and $\bth$ is converging weakly, this lemma implies that $\|\bth\|_{\ell_k}^k$ is bounded wpa.~1.}
\end{proof}



\cmt{
\[
\te{\tn{\Vb\eb}^2-\E[\tn{\Vb\eb}^2]}\lesssim 1/\sqrt{p}.
\]
}

\cmt{
Set $\A= (\Qb\bSi\Qb^T)^{-1/2}$. Observe that $\A\lesssim \Iden_n$ and that $\Vb=\A\Qb\sqrt{\bSi}$.
Set $\bar{X}=\eb^T\Qb^T\Qb\bts$. Observing $\Qb\bSi\Qb^T\gtrsim \Sigma_{\min}\Iden_{n}$, this yields that
\[
\E|X|^k\lesssim 
\]
Note that $\E[X]\lesssim 1/\sqrt{p}$ from above. Secondly, }









\cmt{
Let $\Qb$ be a uniformly random subspace satisfying $\Qb\Qb^T=\Iden_n$. Writing $\X=\bar{\X}\sqrt{\bSi}$ and noticing row space of $\bar{\X}$ is statistically identical to $\Qb$,
Now, let
%
\[
\A=\E[\Qb\bSi\Qb^T]^{-1/2}.
\]
Observe that, upper and lower eigenvalues of $\A$ are bounded by constants. Additionally, it follows from the properties of random tight frames that $\A$ is a diagonal matrix (specifically, using rows of $\Qb$ are symmetrically distributed). Now, observe that
\[
\A\Qb\bSi\Qb^T\A=\Iden.
\] 
Consequently $\A\Qb\sqrt{\bSi}\sim \Vb$. Thus, we can write
\[
\Vb^T\Vb\bt=\sqrt{\bSi}\Qb\A^2\Qb^T\sqrt{\bSi}\bt.
\]
We shall argue that, this vector is subexponential. First, given fixed vector $\vb$, we have that
\[
\E[\Qb\A^2\Qb^T\vb]
\]
we bound the infinity norm of the mean given by%$\|\E[\Vb^T\Vb\bt]\|_{\ell_{\infty}}$. Clearly
\[
\|\E[\Vb^T\Vb\bt]\|_{\ell_{\infty}}\leq 
\] }
%%%%%%%%%%%%%%%%%%%%%%%%%%%%%%%%%


Lemma \ref{lem:bd_PO} implies that nothing changes in \eqref{eq:PO} if we further constrain $\w\in\Bc$ in \eqref{eq:PO}. Henceforth, with some abuse of notation, we let
\begin{align}\label{eq:PO_bd}
\Phi(\X):=\min_{\w\in\Bc} \frac{1}{2}\tn{\Sigmab^{-1/2}\w+\betas}^2\quad\text{subject to}\quad \Xb\w=\sigma \z,
\end{align}

Next, in order to analyze the primary optimization (PO) problem in \eqref{eq:PO_bd} in apply the CGMT \cite{thrampoulidis2015regularized,thrampoulidis2018precise}. Specifically, we use the constrained formulation of the CGMT given by Theorem \ref{thm closed}.  Specifically, the auxiliary problem (AO) corresponding to \eqref{eq:PO_bd} takes the following form with $\g\sim\Nn(0,\Iden_n)$, $\h\sim\Nn(0,\Iden_p)$, $h\sim \Nn(0,1)$:% (non-rigorously)
\begin{align}
\phi(\g,\h) = \min_{\w{\in\Bc}} \frac{1}{2}\tn{\Sigmab^{-1/2}\w+\betas}^2\quad\text{subject to}\quad \tn{\g}\tn{\w~{\sigma}}\leq \h^T\w+\sigma h.\label{eq:AO_con}
\end{align}%\tn{\h}
%Set $\bar{\h}=\h/\sqrt{p}$.

We will prove the following techincal result about the AO problem.
%\somm{Re-specify the assumptions of this lemma!}
\begin{lemma}[Properties of the AO -- Overparameterized regime]\label{lem:AO}
{Let the assumptions of Theorem \ref{thm:master_W2} hold.} Let $\phi_n=\phi(\g,\h)$ be the optimal cost of the minimization in \eqref{eq:AO_con}. Define $\bar\phi$ as the optimal cost of the following deterministic min-max problem
\begin{align}\label{eq:AO_det}
\bar\phi:=\max_{u\geq 0}\min_{\tau>0}~ \Dc(u,\tau):=\frac{1}{2}\left({u\tau} + \frac{u\sigma^2}{\tau} - u^2\kappa\,\E\left[ \frac{1}{\Lambda^{-1}+\frac{u}{\tau}} \right] - {\E\left[\frac{B^2}{1+\frac{u}{\tau}\Lambda}\right]} \right).
\end{align}
The following statements are true.
%\somm{Important: Added a AO subscript to AO related terms!}

\noindent{(i).}~The AO minimization in \eqref{eq:AO_con} is $\frac{1}{\Sigma_{\max}}$-strongly convex and has a unique minimizer $\wha_n:=\wha_n(\g,\h)$.

\noindent{(ii).}~In the limit of $n,p\rightarrow\infty, p/n=\kappa$, it holds that $\phi(\g,\h)\rP\bar\phi$, i.e., for any $\eps>0$:
$$
\lim_{n\rightarrow\infty}\P\left(|\phi(\g,\h)-\bar\phi|>\eps\right) = 0.
$$

\noindent{(iii).} The max-min optimization in \eqref{eq:AO_det} has a unique saddle point $(u_*,\tau_*)$ satisfying the following: 
$$
u_*/\tau_* = \xi\quad\text{and}\quad\tau_* = \gamma,
$$
where $\xi, \gamma$ are defined in Definition \ref{def:Xi}.


\noindent{(iv).}~Let $f:\R^3\rightarrow\R$ be a function in \fx{$\rm{PL}(3)$}. Let $\btha_n=\btha_n(\g,\h)=\Sigmab^{-1/2}\wha_n + \betas$. Then,
$$
\frac{1}{p}\sum_{i=1}^{p}f\left(\sqrt{p}\btha_n,\sqrt{p}\betas,\Sigmab\right) \rP \E_{(B,\Lambda,H)\sim\mu\otimes	\Nn(0,1)}\left[f\left(X_{\kappa,\sigma^2}(B,\Lambda,H),B,\Lambda\right) \right].
$$
\fx{In particular, this holds for all functions  $f\in \Plt$ defined in \eqref{eq:pdef}.}%\CT{Stating this way so that it becomes more clear where the $W_4$ convergence requirement comes from. See also remark below.}


\noindent{(v).}~\cts{The empirical distribution of $\btha_n$ converges weakly to the measure of $X_{\kappa,\sigma^2}$, and also,  for some absolute constant $C>0$:
\begin{align}\label{eq:k_AO}
\tn{\btha}^2 < C\quad\text{wpa 1.}
%\sum_{i=1}^{p}(\btha_n)^2
% \rP \E\left[(X_{\kappa,\sigma^2})^{2k-2}\right] < \infty.
\end{align}
}



\end{lemma}




%\begin{lemma}\label{lem:deviation}
%Let $\wh_n=\wh_n(\X,\z)$ be a minimizer of the PO in \eqref{eq:PO}. 
%%Further let $\tau_n=\tau_n(\g,\h,\z)$ and $u_n=u_n(\g,\h,\z)$ be the unique saddle point of the min-max optimization in \eqref{eq:AO_4}. Define $\wh_n=\wh_n(\tau_n,u_n)$ as in \eqref{eq:w_n}. 
%For any pseudo-Lipschitz function $f:\R\rightarrow\R$ of order 2, it hold that
%\begin{align}
%\frac{1}{p} \sum_{i=1}^{p} f\left(\hat\w_{n,i}\right) \rP \E_\mu\left[X_{\kappa,\sigma^2}(\Lambda,N,H)  \right].
%\end{align}
%\end{lemma}


%\begin{proof}[Proof of Lemma \ref{lem:deviation}]
We prove Lemma \ref{lem:AO} in Section \ref{sec:proofAO}. \fx{We remark that Assumption \ref{ass:mu} on $W_4$-convergence of the joint empirical distribution of $\{(\Sigmab_{i,i},\sqrt{p}\bts_i)\}_{i\in[p]}$ is required in the proof of the statement (iv) above. More generally if $W_k$-convergence is known for some integer $k$, then statement (iv) above holds for test functions $f\in\rm{PL}(k-1)$. This is the first place in the proof of Theorem \ref{thm:master_W2}, where we use the assumption $f\in\Plt$; indeed, we show in Lemma \ref{lem:fL} that $\Fl\subset\Plt\subset\rm{PL}(3)$. The second part is in proving the perturbation result in \eqref{eq:dev2show} below. Unlike the former, when proving the perturbation result, the requirement $f\in\Plt$ cannot be relaxed (e.g.~to $f\in\rm{PL}(k-1)$) by simply increasing the order of $W_k$-convergence in Assumption \ref{ass:mu}.}


\subsubsection{Finalizing the proof of Theorem \ref{thm:master_W2}:}
Here, we show how Lemma \ref{lem:AO} leads to the proof of Theorem \ref{thm:master_W2} when combined with the CGMT framework \cite{thrampoulidis2015regularized,thrampoulidis2018precise}.

\cts{Let $f:\R^3\rightarrow\R$ be a function in $\Plt$, where $\Plt$ was defined in \eqref{eq:pdef}.} For convenience, 
 define 
$$F_n(\bth_n,\betas,\Sigmab):=\frac{1}{p} \sum_{i=1}^{p} f\left(\sqrt{p}\bth_{n,i},\sqrt{p}\betas_{i},\Sigmab_{ii}\right)\quad\text{and}\quad\alpha_*:=\E_\mu\left[f(X_{\kappa,\sigma^2}(\Lambda,B,H),B,\Lambda)\right].$$
Fix  any $\eps>0$ and define the set 
\begin{align}\label{eq:S_set}
\Sc = \Sc(\betas,\Sigmab) 
%=\{\bt\bgl |F_n(\bth_n,\betas,\Sigmab)-\alpha_*|\geq 2\eps\}.
=\big\{\w=\sqrt{\Sigmab}(\betab-\bts) \in\Bc \bgl |F_n(\bt,\betas,\Sigmab)-\alpha_*|\geq 2\eps\big\}.
\end{align}
\fy{With this definition, observe that, it suffices to prove that the solution $\hat\w_n=\sqrt{\Sigmab}(\bth_n-\bts)$ of the PO in \eqref{eq:PO_beta} satisfies $\hat\w_n\not\in\Sc$ wpa 1.}
% To see that this is sufficient, note the following. On the one hand, setting $f=\fl\in \Fl$, directly proves \eqref{eq:thm}. On the other hand, recall that $W_2$-convergence is equivalent to convergence of any $\rm{PL}(2)$ test function $f$, e.g. see \cite[Sec. 6.1]{montanari2017estimation}. 

%\ct{Note that this $\Sc$ is a random set now.}
To prove the desired, we need to consider the ``perturbed" PO and AO problems (compare to \eqref{eq:PO} and \eqref{eq:AO_con}) as:
\begin{align}\label{eq:PO_S}
\Phi_S(\X)=\min_{\w\in\Sc} \frac{1}{2}\tn{\Sigmab^{-1/2}\w+\betas}^2\quad\text{subject to}\quad \Xb\w=\sigma \z,
\end{align}
and
\begin{align}
\phi_S(\g,\h)=\min_{\w\in\Sc} \frac{1}{2}\tn{\Sigmab^{-1/2}\w+\betas}^2\quad\text{subject to}\quad \tn{\g}\tn{\w~{\sigma}}\leq \h^T\w+\sigma h.\label{eq:AO_S}
\end{align}%\tn{\h}
Recall here, that $\Xb=\X\Sigmab^{-1/2}\distas\Nn(0,1)$, $\g\sim\Nn(0,\Iden_n)$, $\h\sim\Nn(0,\Iden_p)$, $h\sim \Nn(0,1)$ and we have used the change of variables $\w:=\sqrt{\Sigmab}(\bt-\betas)$ for convenience. 

 Using \cite[Theorem 6.1(iii)]{thrampoulidis2018precise} it suffices to find costants $\bar\phi, \bar\phi_S$ and $\eta>0$ such that the following three conditions hold:
\begin{enumerate}
\item $\bar\phi_S \geq \bar\phi + 3\eta$,
\item $\phi(\g,\h) \leq \bar\phi + \eta$, with probability approaching 1,
\item $\phi_S(\g,\h) \geq \bar\phi_S - \eta$, with probability approaching 1.
\end{enumerate}
In what follows, we  explicitly find $\bar\phi, \bar\phi_S,\eta$ such that the three conditions above hold.

\vspace{5pt}
\noindent\underline{Satisfying Condition 2}: Recall the deterministic min-max optimization in \eqref{eq:AO_det}. Choose $\bar\phi=\Dc(u_*,\tau_*)$ be the optimal cost of this optimization. From Lemma \ref{lem:AO}(ii), $\phi(\g,\h)\rP\bar\phi$. Thus, for any $\eta>0$, with probability approaching 1:
\begin{align}\label{eq:phi_lim}
\bar\phi + \eta \geq \phi(\g,\h) \geq \bar\phi - \eta.
\end{align}
Clearly then, Condition 2 above holds for any $\eta>0$. 

\vspace{5pt}
\noindent\underline{Satisfying Condition 3}: Next, we will show that the third condition holds for appropriate $\bar\phi$. 
%Let $\tau_n=\tau_n(\g,\h,\z)$ and $u_n=u_n(\g,\h,\z)$ be the unique saddle point of the min-max optimization in \eqref{eq:AO_4}. Define 
Let $\wha_n=\wha_n(\g,\h)$ be the unique minimizer of \eqref{eq:AO_con} as per Lemma \ref{lem:AO}(i), i.e., $\frac{1}{2}\tn{\Sigmab^{-1/2}\wha_n+\bt}^2 = \phi(\g,\h)$. Again from Lemma \ref{lem:AO}, the minimization in \eqref{eq:AO_con} is $1/\Sigma_{\max}$-strongly convex in $\w$. Here, $\Sigma_{\max}$ is the upper bound on the eigenvalues of $\Sigmab$ as per Assumption \ref{ass:mu}. Thus, for any $\tilde\eps>0$ and any feasible $\w$ the following holds (deterministically):
\begin{align}\label{eq:sc}
\frac{1}{2}\tn{\Sigmab^{-1/2}\w+\bt}^2 \geq \phi(\g,\h) + \frac{\tilde\epsilon^2}{2{\Sigma_{\max}}},~\text{provided that}~ \|\w-\wha_n\|_2 \geq \tilde\epsilon.
\end{align}

Now, we argue that {wpa 1,}
\begin{align}\label{eq:dev_arg}
%~\bt\in\Sc \text{ implies that }
\text{for all}~\w\in \Sc~\text{it holds that}~\|\w-\wha_n\|_2\geq \tilde\eps,
\end{align}
for an appropriate value of a constant $\tilde\eps>0$.

 Consider any $\w\in\Sc$. 

 First, by definition in \eqref{eq:S_set}, for $\betab=\Sigmab^{-1/2}\w+\bts$ we have that
$$
|F_n(\bt,\betas,\Sigmab)-\alpha_*| \geq 2\eps.
$$

Second, by Lemma \ref{lem:AO}(iv), with probability approaching 1,
$$
|F(\btha_n,\betas,\Sigmab) - \alpha_*| \leq \epsilon.
$$

Third, we will show that wpa 1, there exists universal constant $C>0$ such that
\begin{align}
|F_n(\btha_n,\betas,\Sigmab) - F_n(\bt,\betas,\Sigmab)| \leq C {\|\btha_n - \bt\|_2}\label{eq:dev2show}.
\end{align}

Before proving \eqref{eq:dev2show}, let us argue how combining the above three displays shows the desired. Indeed, in that case, wpa 1,
\begin{align*}
2\eps &\leq |F_n(\bt,\betas,\Sigmab)-\alpha_*| \leq |F_n(\btha_n,\betas,\Sigmab) - F_n(\bt,\betas,\Sigmab)| + |F_n(\btha_n,\betas,\Sigmab) - \alpha_*| \\
&\leq \epsilon + C \,\|\bt-\btha_n\|_2. \\
&\qquad\Longrightarrow \|\bt-\btha_n\|_2 \geq {\eps}/{C}=:\hat\eps\\
&\qquad\Longrightarrow \|\w-\wha_n\|_2 \geq \hat\eps\sqrt{\Sigma_{\min}}=:\tilde\eps.
\end{align*}
In the last line above, we recalled that $\bt=\Sigmab^{-1/2}\w+\betas$ and $\Sigmab_{i,i}\geq\Sigma_{\min},~i\in[p]$ by Assumption \ref{ass:mu}. This proves \eqref{eq:dev_arg}.

Next, combining \eqref{eq:dev_arg} and \eqref{eq:sc}, we find that wpa 1,
$
\frac{1}{2}\tn{\Sigmab^{-1/2}\w+\bt}^2 \geq \phi(\g,\h) + \frac{\tilde\epsilon^2}{2\Sigma_{\max}},~\text{for all}~ \w\in\Sc.
$
Thus, 
\begin{align}
\phi_S(\g,\h) \geq \phi(\g,\h) + \frac{\tilde\epsilon^2}{2\Sigma_{\max}}.\nn
\end{align}
%\end{proof}
When combined with \eqref{eq:phi_lim}, this shows that
\begin{align}
\phi_S(\g,\h) \geq \bar\phi + \frac{\tilde\epsilon^2}{2\Sigma_{\max}} - \eta.
\end{align}
Thus, choosing $\bar\phi_S = \bar\phi + \frac{\tilde\epsilon^2}{2\Sigma_{\max}}$ proves the Condition 3 above. 


% by the pseudo-Lipschitz property, boundedness of $\bth_n$ and $\bt$ (see Lemma \ref{lem:bd_PO}) the following holds. 

%%%%%%%%%%%%%%%%%%%%%%%%%%%%%%%%%%%
%%%%%%%%%%%%%.  PL(k) argument!!
%%%%%%%%%%%%%%%%%%%%%%%%%%%%%%%%%%%%%
%To simplify notation, let $\ab_i=(\sqrt{p}\bth_{n,i}(\g,\h),\sqrt{p}\betas_i,\Sigmab_{i,i})$ and $\bb_i=(\sqrt{p}\bt_i,\sqrt{p}\betas_i,\Sigmab_{i,i})$, for $i\in[p]$. Then, using $f\in\rm{PL}(k)$,
%\begin{align}
%|F_n(\bth_n(\g,\h),\betas,\Sigmab) - F_n(\bt,\betas,\Sigmab)| &\leq \frac{L}{p}\sum_{i=1}^p (1+\max\{\|\ab_i\|_2^{k-1},\|\bb_i\|_2^{k-1}\})\|\ab_i-\bb_i\|_2\nn \\
%&= \frac{L}{p}\sum_{i=1}^p \left(1+\max\{\|\ab_i\|_2^{k-1},\|\bb_i\|_2^{k-1}\}\right)\cdot|\sqrt{p}\bth_{n,i}(\g,\h) - \sqrt{p}\bt_i|\nn\\
%&\leq L\left(1+\max\{\frac{1}{p}\sum_{i=1}^p \|\ab_i\|_2^{2k-2},\frac{1}{p}\sum_{i=1}^p\|\bb_i\|_2^{2k-2}\} \right)^{1/2}{\|\bth_n(\g,\h) - \bt\|_2}\label{eq:LipC}\,.
%\end{align}
%The last inequality above follows by applying Cauchy-Schwartz.  
%To reach \eqref{eq:dev2show} from the above, it suffices to show that
%\begin{align}
%\max\{\frac{1}{p}\sum_{i=1}^p \|\ab_i\|_2^{2k-2},\frac{1}{p}\sum_{i=1}^p\|\bb_i\|_2^{2k-2}\} \leq C,\label{eq:leqC}
%\end{align}
%for some universal constant $C$ (that may depend on $k$, but not on $n,p$).
%% $\frac{1}{p}\sum_{i=1}^p \|\ab_i\|_2^{2k-2}<\infty$ and $\frac{1}{p}\sum_{i=1}^p \|\bb_i\|_2^{2k-2}<\infty$ as $p\rightarrow\infty$. But, 
%To prove \eqref{eq:leqC}, note that 
% using $a^{\ell_1}b^{\ell_2}c^{\ell_3}\leq a^{\ell}+b^{\ell}+c^\ell,~\ell=\ell_1+\ell_2+\ell_3$, for some constant $C=C(k)>0$ it holds:
%%\ct{Why is this (or something like this true)?}
%%\som{arithmetic - geometric mean inequality?}
%\begin{align}
%\frac{1}{p}\sum_{i=1}^p \|\bb_i\|_2^{2k-2} {\leq} C\left(\frac{1}{p}\sum_{i=1}^p |\sqrt{p}\bt_{i}|^{2k-2} + \frac{1}{p}\sum_{i=1}^p |\sqrt{p}\betas_{i}|^{2k-2} + \frac{1}{p}\sum_{i=1}^p |\Sigmab_{i,i}|^{2k-2} \right)\label{eq:bdmom}\,.
%\end{align}
%\cts{The terms $\frac{1}{p}\sum_{i=1}^p |\sqrt{p}\betas_{i}|^{2k-2} $ and $ \frac{1}{p}\sum_{i=1}^p |\Sigmab_{i,i}|^{2k-2}$ are bounded by Assumption \ref{ass:mu}.}
%%\ct{Need to add assumption that $ \frac{1}{p}\sum_{i=1}^p |\Sigmab_{i,i}|^{2k-2}<\infty$ and $\frac{1}{p}\sum_{i=1}^p (\sqrt{p}|\betas_{i}|^{2k-2} $. OK by just assuming that $B$ is Definition \ref{def:Xi} is bounded.}. 
%{\color{red}Furthermore we additionally know from Lemma \ref{subexp bth} that, 
%\begin{align}
%\frac{1}{p}\sum_{i=1}^p |\sqrt{p}\bt_{i}|^{2k-2} <\infty
%\end{align}
%}
%%\frac{1}{d}\sum_{i=1}^p |\bt_{i}|^{2k-2}\leq \|\bt\|_2^{2k-2}<C'(k)$, where the last inequality follows from feasibility of $\bt$, i.e., $\w=\sqrt{\Sigmab}(\bt-\betas)\in\Bc$; see Lemma \ref{lem:bd_PO} and \eqref{eq:bd_beta}. 
%These together with \eqref{eq:bdmom} show that $\frac{1}{p}\sum_{i=1}^p \|\bb_i\|_2^{2k-2}<C$ wpa 1. Similarly, we can argue for $\frac{1}{p}\sum_{i=1}^p \|\ab_i\|_2^{2k-2}$, completing the proof of \eqref{eq:leqC}.
%%%%%%%%%%%%%%%%%%%%%%%%%%
%%%%%%%%%%%%%%%%%%%%%%%%%%%%%%%%

\vp
\noindent{\textbf{Perturbation analysis via Pseudo-Lipschitzness (Proof of \eqref{eq:dev2show}).}} To complete the proof, let us now show \eqref{eq:dev2show}. Henceforth, $C$ is used to denote a universal constant whose value can change from line to line. \fy{Recall that $f\in \Plt$ where $\Plt:\R^2\times \Zc\rightarrow\R$ is the set of $\rm{PL}(3)$ functions such that $f(\cdot,\cdot,z)$ is $\rm{PL}(2)$ for all $z\in\Zc$. Suppose that the $\rm{PL}(2)$ constant of $f(\cdot,\cdot,z)$ is upper bounded over $z\in \Zc$ by some $C>0$. We also let $C$ change from line to line for notational simplicity. Then, we have the following chain of inequalities:
\begin{align}
|F_n(\btha_n(\g,\h),\betas,\Sigmab) - F_n(\bt,\betas,\Sigmab)|
&=\frac{1}{p}\sum_{i=1}^p|f(\sqrt{p}\btha_{n,i},\sqrt{p}\betas_i,\Sigmab_{i,i})-f(\sqrt{p}\bt_i,\sqrt{p}\betas_i,\Sigmab_{i,i})|
\nn\\
&\leq
 \frac{C}{p}\sum_{i=1}^p (1+ \|\sqrt{p}[\betas_i,\btha_{n,i}]\|_2 +  \|\sqrt{p}[\betas_i,\bt_{i}]\|_2) \sqrt{p}|\btha_{n,i}-\bt_{i}|\nn\\
&\leq
C \Big(1+ \frac{1}{\sqrt{p}}\big(\sum_{i=1}^{p}\|\sqrt{p}[\betas_i,\btha_{n,i}]\|_2^2\big)^{1/2} + \frac{1}{\sqrt{p}}\big(\sum_{i=1}^p\|\sqrt{p}[\betas_i,\bt_{i}]\|_2^2\big)^{1/2}\Big) \|\btha_n-\bt\|_2\nn\\
&\leq
%C  \left(1+ \|\betas\|_2^2+ \|\btha_n\|_2^2 + \|\bt\|_2^2\right) \|\btha_n-\bt\|_2.
C \left(1+ \max\{\|\betas\|_2^2,\|\btha_n\|_2^2,\|\bt\|_2^2\}^{1/2} \right) \|\btha_n-\bt\|_2.\label{eq:fcase_main}
 \end{align}
 %boundedness of $\lai$ as per Assumption \ref{ass:mu}. {In the third line, we used the fact that the function $\psi(a,b) = (a-g(b))^2$ is $\rm{PL}(2)$ as the quadratic of a Lipschitz function.}
 In the second line above, we used the fact that $f(\cdot,\cdot,z)$ is $\rm{PL}(2)$. The third line follows by Cauchy-Schwartz  inequality. Finally, in the last line, we used the elementary fact that $a+b+c\leq 3\max\{a,b,c\}$ for $a=2\sum_{i=1}^p(\betas_i)^2$ and $b=\sum_{i=1}^p(\btha_{n,i})^2$ and $c=\sum_{i=1}^p\bt_{i}^2$.}

\begin{comment}
 First, consider the case $f=\fl\in\Fl$, i.e., $f(x,y,z) = z(y-g(x))^2$ for Lipschitz $g$. We have the following chain of inequalities:
%that we care about for evaluating the risk (see Lemma 5) we can write \eqref{eq:LipC} simpler as follows:
\begin{align}
|F_n(\btha_n(\g,\h),\betas,\Sigmab) - F_n(\bt,\betas,\Sigmab)| &=
 \frac{1}{p}\sum_{i=1}^p|\lai|\left| (\sqrt{p}\betas_i-g(\sqrt{p}\btha_{n,i}))^2-(\sqrt{p}\betas_i-g(\sqrt{p}\bt_{i}))^2 \right|\nn\\
 &\leq
\Sigma_{\max} \frac{1}{p}\sum_{i=1}^p \left| (\sqrt{p}\betas_i-g(\sqrt{p}\btha_{n,i}))^2-(\sqrt{p}\betas_i-g(\sqrt{p}\bt_{i}))^2 \right|\nn\\
 &\leq
{C} \Sigma_{\max} \frac{1}{p}\sum_{i=1}^p (1+ \|\sqrt{p}[\betas_i,\btha_{n,i}]\|_2 +  \|\sqrt{p}[\betas_i,\bt_{i}]\|_2) \sqrt{p}|\btha_{n,i}-\bt_{i}|\nn\\
 &\leq
{C} \Sigma_{\max}  \Big(1+ \frac{1}{\sqrt{p}}\big(\sum_{i=1}^{p}\|\sqrt{p}[\betas_i,\btha_{n,i}]\|_2^2\big)^{1/2} + \frac{1}{\sqrt{p}}\big(\sum_{i=1}^p\|\sqrt{p}[\betas_i,\bt_{i}]\|_2^2\big)^{1/2}\Big) \|\btha_n-\bt\|_2\nn\\
 &\leq
%C  \left(1+ \|\betas\|_2^2+ \|\btha_n\|_2^2 + \|\bt\|_2^2\right) \|\btha_n-\bt\|_2.
C  \left(1+ \max\{\|\betas\|_2^2,\|\btha_n\|_2^2,\|\bt\|_2^2\}^{1/2} \right) \|\btha_n-\bt\|_2.\label{eq:fcase1}
 \end{align}
 In the second line above, we used boundedness of $\lai$ as per Assumption \ref{ass:mu}. {In the third line, we used the fact that the function $\psi(a,b) = (a-g(b))^2$ is $\rm{PL}(2)$ as the quadratic of a Lipschitz function.} The fourth line follows by Cauchy-Schwartz  inequality. Finally, in the last line, we used the elementary fact that $a+b+c\leq 3\max\{a,b,c\}$ for $a=2\sum_{i=1}^p(\betas_i)^2$ and $b=\sum_{i=1}^p(\btha_{n,i})^2$ and $c=\sum_{i=1}^p\bt_{i}^2$.
 
Second, consider the case $f\in\rm{PL}(2)$. Let $\ab_i=(\sqrt{p}\btha_{n,i},\sqrt{p}\betas_i,\Sigmab_{i,i})$ and $\bb_i=(\sqrt{p}\bt_i,\sqrt{p}\betas_i,\Sigmab_{i,i})$, for $i\in[p]$. A similar chain of inequality holds:
%\somm{The solution for PL(k) is showing that around $\hat{\bt}$, there is an arbitrarily close $\hat{\bt}_{bdd}$ with bounded entries.}\CT{Not sure what this comment refers to.}
\begin{align}
|F_n(\btha_n,\betas,\Sigmab) - F_n(\bt,\betas,\Sigmab)| &\leq \frac{{C}}{p}\sum_{i=1}^p (1+\max\{\|\ab_i\|_2,\|\bb_i\|_2\})\|\ab_i-\bb_i\|_2\nn \\
&= \frac{{C}}{p}\sum_{i=1}^p \left(1+\max\{\|\ab_i\|_2,\|\bb_i\|_2\}\right)\cdot|\sqrt{p}\btha_{n,i} - \sqrt{p}\bt_i|\nn\\
&\leq {C}\left(1+\max\{\frac{1}{p}\sum_{i=1}^p \|\ab_i\|_2^{2},\frac{1}{p}\sum_{i=1}^p\|\bb_i\|_2^{2}\} \right)^{1/2}{\|\btha_n - \bt\|_2}\nn\\
&\leq C\left(1+\max\{\|\btha_n\|_2^2,\|\betas\|_2^2,\|\bt\|_2^2,\frac{1}{p}\sum_{i=1}^p\lai^2\} \right)^{1/2}{\|\btha_n - \bt\|_2}\nn\\
&\leq C\Sigma_{\max}^2\left(1+\max\{\|\btha_n\|_2^2,\|\betas\|_2^2,\|\bt\|_2^2\} \right)^{1/2}{\|\btha_n - \bt\|_2}.\label{eq:fcase2}
\end{align}
The first inequality uses the fact that $f\in\rm{PL}(2)$. The second inequality in the third line follows by Cauchy-Schwartz. The last inequality used Assumption \ref{ass:mu} on boundedness of $\lai$.
\end{comment}

Hence, it follows from \eqref{eq:fcase_main} that in order to prove \eqref{eq:dev2show}, we need to show boundedness of the following terms: $\|\btha_n\|_2$, $\|\betas\|_2$ and $\|\bt\|_2$. {By feasibility of $\btha_n$ and $\bt$, we know that $\btha_n,\bt\in\Bc$. Thus, the desired $\|\bt\|_2<\infty$ and $\|\btha_n\|_2<\infty$ follow directly by Lemma \ref{lem:bd_PO} (Alternatively, for $\btha_n$ we conclude the desired by directly applying Lemma \ref{lem:AO}(v)).} Finally, to prove $\|\betas\|_2<\infty$, note that
$$
\|\betas\|_2^2 =\frac{1}{p}\sum_{i=1}^p(\sqrt{p}\betas_i)^2,
$$
\cts{which is bounded wpa 1 by Assumption \ref{ass:mu}, which implies bounded second moments of $\sqrt{p}\betas$. } This completes the proof of \eqref{eq:dev2show}, as desired.
%And everything is easy now since $\|\betas\|_2^2$ is bounded as long as second moments of empirical distribution as bounded, $\|\bth_n\|_2^2$ is bounded if $\|\betas\|_2^2$ is bounded, and, $\|\bt\|_2^2$ is bounded, since $\bt\in\Bc$ by Lemma 1. To prove 


\vspace{5pt}
\noindent\underline{Satisfying Condition 1:} To prove Condition 1, we simply pick $\eta$ to satisfy the following
\begin{align}
\bar\phi_S > \bar\phi + 3 \eta ~\Leftarrow~ \bar\phi + \frac{\tilde\epsilon^2}{2\Sigma_{\max}} - \eta \geq \bar\phi + 3 \eta ~\Leftarrow~ \eta \leq \frac{\tilde\epsilon^2}{8\Sigma_{\max}}.\nn
\end{align}


This completes the proof of Theorem \ref{thm:master_W2}.

~~~~
~~~~
%%
%%
%%\begin{lemma} The function $h$ in \eqref{h func} is $\rm{PL}(k)$ as long as $f$ is $\rm{PL}(k)$.
%%{\color{red}The function $h$ is $\R^3\rightarrow\R$, not $\R^{3p}\rightarrow\R$! See Lemma 6 below.}
%%\end{lemma}
%%\begin{proof} 
%%Since $f$ is $\rm{PL}(k)$, for some $L>0$, \eqref{PL func} holds. Fix $\ab=[\x,\y,\z]\in\R^{3p},\ab'=[\x',\y',\z']\in\R^{3p}$. Let $\bb=[\g,\y,\z]$ where $\g=g(\y,\z,\x)$. We have that
%%\begin{align}
%%|h(\ab)-h(\ab')|&\leq |f(\bb)-f(\bb')|\\
%%&\leq L\left(1+\|\bb\|_2^{k-1}+\|\bb'\|_2^{k-1}\right)\|\bb-\bb'\|_2\\
%%&\lesssim L\left(1+\|\ab\|_2^{k-1}+\|\ab'\|_2^{k-1}+\|\g\|_2^{k-1}+\|\g'\|_2^{k-1}\right)(\|\ab-\ab'\|_2+\|\g-\g'\|_2).
%%%~~~~~~~~&L\left(1+\|\g(\y,\z,\x)\|_2^{k-1}+\|g(\y',\z',\x')\|_2^{k-1}\right)\|g(\y,\z,\x)-g(\y',\z',\x')\|_2.
%%\end{align}
%%Next, we need to bound the $\g$ term in terms of $\ab$. This is accomplished as follows
%%\begin{align}
%%\|\g\|_2^{k-1}&\lesssim \frac{1}{p}\sum_{i=1}^p|g_i|^{k-1}\\
%%&\lesssim \frac{1}{p}\sum_{i=1}^p\left|\frac{y_i^{-1/2}}{1+(\ksi y_i)^{-1}}\sqrt{\kappa}\tau_* z_i + (1-(1+\ksi_*y_i)^{-1})x_i\right|\\
%%&\lesssim \frac{1}{p}\sum_{i=1}^p(|z_i|+|x_i|)^{k-1}\lesssim \|\x\|_2^{k-1}+\|\z\|_2^{k-1}\\
%%&\lesssim \|\ab\|^{k-1}.
%%\end{align}
%%Secondly and similarly, we have the following perturbation bound on the $\g-\g'$. \so{Suppose the triples $(x_i,y_i,z_i)$ takes values in a fixed bounded compact set $\Mc$. We will prove the following sequence of inequalities} 
%%\begin{align}
%%%|\|\g\|_2-\|\g'\|_2|&\leq 
%%\tn{\g-\g'}^2&\leq \sum_{i=1}^p|g(x_i,y_i,z_i)-g(x'_i,y'_i,z'_i)|^2\\
%%&\leq C^2_{\Mc} \sum_{i=1}^p(|x_i-x'_i|^2+|y_i-y'_i|^2+|z_i-z'_i|^2)\label{sec line eq}\\
%%&\leq C^2_{\Mc} (\tn{\x-\x'}^2+\tn{\y-\y'}^2+\tn{\z-\z'}^2)\\
%%&\lesssim \tn{\ab-\ab'}^2.
%%\end{align}
%%In what follows, we prove the second line \eqref{sec line eq} i.e.~the fact that for any triples $a=(x,y,z),a'=(x',y',z')$ (with $a,a'\in\R^3$), we have that 
%%\[
%%|g(a)-g(a')|\leq C_{\Mc}\tn{a-a'}.
%%\]
%%Set $C_\nabla=\sup_{a\in \Mc}\tn{\nabla g(a)}$. By definition of gradient, the inequality above holds with $C_\Mc=C_\nabla$. Thus, all that remains is proving that $C_\nabla$ is upper bounded by a constant. \so{However, this automatically holds because from the definition of $g(\dots)$ function, it is clear that $C_\nabla$ is defined everywhere and continuous thus it has a finite maximum over a compact set.} \cts{Only the problem is that $z_i$ (i.e., $\h_i$) is $\Nn(0,1)$ so it is not bounded.}
%%% is bounded by a constant.
%%\end{proof}
%%


%But from Assumption \ref{ass:mu} and $f\in\rm{PL(2)}$ it follows immediately that with probability approaching 1,
%\begin{align}
%|\frac{1}{p}\sum_{i=1}^{p}f(\sqrt{p}\vb_{n,i})-\E_\mu\left[X_{\kappa,\sigma^2}(\Lambda,N,H)  \right]| \leq \eps.
%\end{align}
%Combining the above displays and using triangle inequality completes the proof of the statement.





















